\begin{abstract}
In this paper, we introduce and study the properties of $\omega$-recursive sequence spaces, a new class of mathematical objects that bridge recursive function theory and functional analysis. For a given recursive function $f: \mathbb{N} \rightarrow \mathbb{N}$, we define the $\omega$-recursive sequence space $\mathcal{R}_f$ as the set of all sequences $(x_n)_{n \in \mathbb{N}}$ satisfying the condition:

\[\sum_{n=1}^{\infty} |x_n|^{f(n)} \leq 1\]

Our main result (Theorem 1) establishes that for any primitive recursive function $f$, the space $\mathcal{R}_f$ is complete under a naturally defined metric and exhibits interesting topological properties related to computational complexity. Through a series of lemmas, we demonstrate that these spaces form a proper hierarchy based on the growth rate of $f$, and that they maintain certain invariance properties under recursive transformations. This work establishes a novel connection between classical sequence spaces and computability theory, opening new avenues for investigating the relationship between analytical and computational properties of infinite sequences.
\end{abstract}